\vspace{20pt}
\section{电容与电容器}

\subsection{电容}
电荷在导体表面的分布必须满足导体静电平衡的条件，对于大小形状确定的孤立导体，由于没有外界电场的影响，根据\Refx{定律}{静电场的唯一性原理}，当导体的所带电荷量确定时，导体外的电场就确定了，导体的电势也就确定了。更进一步的说，由于能使得导体内部场强为零的电荷分布方式是唯一的，当导体的电荷量增加时，并不会改变电荷分布方式，只会等比例的增大各处的电荷，因此，\textbf{导体的电荷量增加若干倍时，导体的电势也将增加若干倍}。
\begin{BoxPrinciple}[孤立导体的电容特性]
    孤立导体的电荷量与其电势成正比。
\end{BoxPrinciple}
定义孤立导体的电荷量与电势的比值，为该孤立导体的\uwave{电容}，通常用$C$表示
\begin{BoxDefinition}[孤立导体的电容]
    C=\frac{q}{V}
\end{BoxDefinition}
\Refx{定义}{孤立导体的电容}中，孤立导体的电容由导体的形状大小决定，对于一般的孤立导体，其电容的计算往往是复杂的，作为一个简单的示例，我们考虑孤立导体球。

根据\Refx{公式}{孤立导体球的电势}和\Refx{定义}{孤立导体的电容}
\begin{BoxEquation}[孤立导体球的电容]
    C=\frac{4\pi\varepsilon_0}{R}
\end{BoxEquation}
由此可见，对于孤立导体球而言，导体球的半径越大，导体球的电容就越大。

\Refx{定义}{孤立导体的电容}中，电容是导体升高单位电势时电荷量的增加量，电容不是导体存储电荷的总容量，这是没有上限的。如果将导体想象为某种存储电荷的仓库，那么，电容并非是仓库的总容量，电容是单位面积的仓库容量，电势就是仓库已使用的面积。
\begin{itemize}
    \item 电容较大，仓库单位面积就可以存储更多电荷，或者说当我们向仓库存入一定量的电荷时，只需要占用较小的仓库面积，即，电容较大的导体，单位电势增量对应较大的电荷增量，反过来说，导体增加一定量的电荷时，导体电势增量较小。
    \item 电容较小，仓库单位面积就只能存储更少电荷，或者说当我们向仓库存入一定量的电荷时，就需要占用较大的仓库面积，即，电容较小的导体，单位电势增量对应较小的电荷增量，反过来说，导体增加一定量的电荷时，导体电势增量较大。
\end{itemize}
由此可见，电容表示导体单位电势存储的电荷的多寡，即导体对电荷的存储效率。

电容在国际单位制中的单位是\si{F=C\cdot V^{-1}}，称为\uwave{法拉}，量纲是\si{[L^{-1}M^{-1}T^4A^2]}。

\subsection{电容器}
孤立导体的电荷量与电势成正比，即$q=CV$，但当导体附近有其他带电导体时，即便导体所带电荷量一定，随着其他导体的位置改变，导体外的电场将发生改变，导体的电势也随之变化，因此非孤立导体的电荷量与电势间并不存在一一对应的正比关系。\vspace{5pt}
\begin{Figure}{电容器}
    \inputfigure[scale=0.75]{Chapter09B_Figure01}
\end{Figure}
如\Refx{图}{电容器}所示，使用接地封闭导体壳$B$将导体$A$包围起来，根据\Refx{定律}{导体空腔和内表面的电荷量关系}，导体$A$带电$q$时，导体壳$B$的内表面也将带电$-q$。

根据\Refx{定律}{空腔接地导体的基本性质}，无论外部电场如何变换，导体$A$与导体壳$B$所带的电荷量与空腔内的电场均不会发生变化，而导体壳$B$接地电势一定，因此导体$A$相对导体壳$B$的电压$U_{AB}$就是一定的，由此导体$A,B$间的电压就与导体$A$的电荷量成正比。

这样一种特殊的导体组称为\uwave{电容器}，而两个导体分别称为电容器的两个\uwave{极板}。

通过以上讨论可知，电容器的两个极板总是在静电感应下带等量异号电荷，且极板上带的电荷量总是与极板间的电压成正比，这一比例系数我们也称为电容
\begin{BoxDefinition}[电容]
    C=\frac{q}{U}
\end{BoxDefinition}
\Refx{定义}{电容}中，应注意符号问题，如果$q$代表导体$A$的电荷量$q_A$，那么$U$就应代表导体$A$至导体壳$B$的电压$U_{AB}$，这样的定义使得电容的值总是为正值。
\begin{itemize}
    \item 当$q_A>0$，此时$q_B<0$，由正电荷至负电荷电势下降，电压为正。
    \item 当$q_A<0$，此时$q_B>0$，由负电荷至正电荷电势上升，电压为负。
\end{itemize}
\Refx{定义}{电容}与先前的\Refx{定义}{孤立导体的电容}是相互兼容的，导体壳$B$接地使得其电势为零，如果我们将导体壳$B$的半径不断扩大直至无穷远处，那么这就相当于孤立导体无穷远处的电势为零，换言之，\textbf{孤立导体是将无穷远处作为另一极板的电容器}，只不过此时由于另一极板的电势恒为零，因此极板间的电压$U$就等同于导体的电势$V$。

实际上，电容器在使用时未必会将外壳接地，电容器的两块极板通常都会接入电路，这并不会改变电容的基本属性，即电容无论如何连接，其电荷与电压成正比的特性不会改变。

下面我们来研究一些特殊形状的电容器的电容大小的计算方法。

\subsubsection{平行板电容器}
我们现在研究\uwave{平行板电容器}，平行板电容器由两块平行放置的金属板构成，这是最常见的一类电容器，也是电容器符号的来源，记金属板的面积为$S$，记金属板的间距为$d$。
\begin{Figure}{平行板电容器}
    \subfigure[平行板电容器]
    {
        \inputfigure[scale=0.75]{Chapter09B_Figure02}
    }\qquad\qquad
    \subfigure[电容器的记号]
    {
        \begin{circuitikz}
            \draw (-2,0) to[C=$C$] (2,0);
        \end{circuitikz}
    }
\end{Figure}
首先解决一个疑问，平行板电容器的结构并不是包围导体的封闭导体壳，并不符合电容器的定义，却仍可以称为电容器？我们知道，封闭曲面的要素包含“将空间划分为两个相互不连通的区域”以及“使得一块区域体积有限，使得另一块区域体积无限”两点，如果我们放弃其中的第二点要素，那么无限大平面在某种意义上也可以视为一种闭曲面。

因此两个平行的无限大带电导体平面也可以构成一个电容器，而当平行板电容器的导体间距远远大于导体面积的线度时，即满足$d\gg\sqrt{S}$时，平行板可以视为是无限大的，平行板电容器也就可以视为电容器，因此两块平行板总是带等量异号的电荷$+Q,-Q$。

现在我们就来计算平行板电容器的电容与其几何形状的关系。

根据\Refx{公式}{等量异号的无限大带电平面间的电场}，平行板电容器间的电场为匀强电场，由于平行板的面积为$S$，若平行板的电荷密度为$\sigma$，那就有$Q=\sigma S$，因此
\begin{equation*}
    E=\frac{\sigma}{\varepsilon_0}=\frac{Q}{\varepsilon_0S}
\end{equation*}
根据\Refx{公式}{匀强电场的电势差}
\begin{equation*}
    U=Ed=\frac{Qd}{\varepsilon_0 S}
\end{equation*}
根据\Refx{定义}{电容}
\begin{equation*}
    C=\frac{Q}{U}=\frac{\varepsilon_0S}{d}
\end{equation*}
即
\begin{BoxEquation}[平行板电容器的电容]
    C=\frac{\varepsilon_0S}{d}
\end{BoxEquation}
\Refx{公式}{平行板电容器的电容}指出，电容正比于极板面积，电容反比于极板间距。较大的电容在电路中是有用途的，而该公式提出了增大电容器的电容的两种途径
\begin{enumerate}
    \item 增大极板面积可以扩大电容，这会导致电容体积的增大。
    \item 减小极板间距可以扩大电容，这看似没有副作用，但事实上会导致电容器耐压的下降。
\end{enumerate}
电容器的耐压是什么？理想的电容器的两块极板间可以承受任意大的电压，但现实的电容器的情况是，极板间的电压升高，极板上的电荷量增大，极板表面的电场强度也随之增大。

如果这种场强大到足矣电离空气分子，这就会使得空气被导通并导致瞬时的剧烈放电，这时我们就称电容器被\uwave{击穿}了，击穿时的巨大能量会导致电容器的结构被破坏，因此每个电容器都具有一个最大可以承受的电压，称为\uwave{耐压}或\uwave{击穿电压}，超过耐压就会使得电容器被击穿。

电容器的耐压与什么因素有关？我们知道，电容器是否会被击穿的关键在于极板的表面场强是否足矣电离空气，这场强记为$E_x$，这场强对于同种材料制成的电容器应为一常数。

如果记电容器的耐压为$V_{\max}$，那么就有
\begin{BoxEquation}[电容器的耐压]
    V_{\max}=E_xd
\end{BoxEquation}
\Refx{公式}{电容器的耐压}指出，电容器的极板面积一定时，电容器的耐压和电容是矛盾的
\begin{itemize}
    \item 增大电容器的耐压，就要增大极板间距$d$，但这会使得电容器的电容下降。
    \item 增大电容器的电容，就要减小极板间距$d$，但这会使得电容器的耐压下降。
\end{itemize}
反过来看，如果我们需要一个具有一定耐压的电容，我们就需要维持一定的极板间距。

假设我们需要一个极板间距为$1$微米，即$d=1\times 10^{-6}$\si{m}的电容，如果希望其具有$1$法拉的电容，根据\Refx{公式}{平行板电容器的电容}，那么极板面积将达到$S=1.13\times 10^{5}$\si{m^2}，这是非常巨大的。因此，法拉是一个很大的单位，法拉级电容不可能通过这种手段制造。

当然法拉级电容是存在的，在下一章中会看到还可以通过填充电介质的方式增加电容。

由于电容器存在击穿电压\footnote{若仍以存储电荷的仓库为比喻，仓库的总面积就是击穿电压，因此，仓库能存储的电荷是有限的（不然会爆仓）。}，因此电容器存储的电荷量也是有上限的，记为$Q_{\max}$
\begin{BoxEquation}[电容器的最大贮电量]
    Q_{\max}=CV_{\max}=E_xS\varepsilon_0
\end{BoxEquation}
由此可见，当电容器极板面积一定时，改变电容器的极板间距，虽然会使得电容器的电容和耐压发生变换，但是并不会改变两者的乘积，即电容器的最大贮电量$Q_{\max}$。

\subsubsection{柱形电容器}
\uwave{柱形电容器}由两个同轴导体圆筒$A,B$作为极板构成，圆筒半径分别为$R_A,R_B$，圆筒径向长度为$L$，通常情况下$L\gg(R_B-R_A)$，因此可以近似将圆筒看作无限长。

设内外圆筒的电荷线密度分别为$\pm\lambda$，根据\Refx{公式}{均匀无限长带电直棒的电场}
\begin{equation*}
    E=\frac{\lambda}{2\pi\varepsilon_0}\frac{1}{r}
\end{equation*}
因此两柱形电极$A,B$间的电压$U=U_{AB}$为
\begin{equation*}
    U=\Int[R_A][R_B]{\frac{\lambda}{2\pi\varepsilon_0}\frac{1}{r}\dif{r}}=\frac{\lambda}{2\pi\varepsilon_0}\qty(\ln R_B-\ln R_A)
\end{equation*}
根据\Refx{定义}{电容}，由于$Q=\lambda L$，柱形电容器的电容为
\begin{BoxEquation}[柱形电容器的电容]
    C=\frac{2\pi\varepsilon_0 L}{\ln R_B-\ln R_A}
\end{BoxEquation}

\subsubsection{球形电容器}
\uwave{球形电容器}由两个同心导体球壳$A,B$作为极板构成，球壳半径分别为$R_A,R_B$。

设内外球壳的电荷量分别为$\pm Q$，根据\Refx{公式}{带电球壳的电场分布}
\begin{equation*}
    E=\frac{Q}{4\pi\varepsilon_0}\frac{1}{r^2}
\end{equation*}
因此两球形电极$A,B$间的电压$U=U_{AB}$为
\begin{equation*}
    U=\Int[R_A][R_B]{\frac{Q}{4\pi\varepsilon_0}\frac{1}{r^2}\dif{r}}=\frac{Q}{4\pi\varepsilon_0}\qty(\frac{1}{R_A}-\frac{1}{R_B})=\frac{Q}{4\pi\varepsilon_0}\frac{R_B-R_A}{R_BR_A}
\end{equation*}
根据\Refx{定义}{电容}，球形电容器的电容为
\begin{BoxEquation}[球形电容器的电容]
    C=\frac{4\pi\varepsilon_0 R_BR_A}{R_B-R_A}
\end{BoxEquation}

\subsection{RC电路}
我们现在来研究电容接入电路时表现出的特性，首先应当说明的是，包含电容的电路与过往高中所熟悉的仅包含电阻的电路是不同的，包括电流电压在内的物理量将不再是常量，而是随时间变化的，或者说是\uwave{时变}的，因此我们解出的电流电压将是一个关于时间的函数。
\begin{Figure}{RC电路}
    \inputfigure{Chapter09B_Figure03}
\end{Figure}
\Refx{图}{RC电路}是一个最为简单的包含电容的电路
\begin{itemize}
    \item 左侧的回路称为~\uwave{RC充电电路}，当$S_1$闭合$S_2$断开时，电容将会充电。
    \item 右侧的回路称为~\uwave{RC放电电路}，当$S_2$闭合$S_1$断开时，电容将会放电。
\end{itemize}
两者统称为~\uwave{RC电路}，即由电阻（\textbf{R}esistor）和电容（\textbf{C}apacitor）构成的电路。

现在我们来具体研究电容的充电过程和放电过程是如何进行的。

\subsubsection{电容的充电过程}
% 下面研究RC充电电路，$S_1$闭合，$S_2$断开。

设电源电动势为$E$，设电阻器的电阻为$R$，设电容器的电容为$C$，这均是常量。

设电流为$I$，设电容器极板间的电压为$U$，设电容器极板上的电荷量为$Q$，这均是函数。

现在我们要求解电容器的电压$U$关于时间的函数，列出以下方程组
\begin{Equation}[电容充电1]
    \begin{cases}
        E=IR+U\\[2mm]
        I=\Fdif{Q}{t}\\[2mm]
        Q=CU
    \end{cases}
\end{Equation}
将式\ref{式:电容充电1}中的(2)式和(3)式先后代入(1)式，消去$I,Q$
\begin{Equation}[电容充电2]
    E=RC\Fdif{U}{t}+U
\end{Equation}
式\ref{式:电容充电2}是一个关于$U$的一阶微分方程。

式\ref{式:电容充电2}可以通过分离变量法求解，分离变量并两侧分别积分得
\begin{Equation}[电容充电3]
    \Int[0][U]{\frac{RC}{E-U}\dif{U}}=
    \Int[0][t]{\dif{t}}
\end{Equation}
现在求解式\ref{式:电容充电3}
\begin{align*}
    -&RC\ln(E-U)|_0^U=t\\[2mm]
    -&RC[\ln(E-U)-\ln E]=t\\[2mm]
    &\ln(\frac{E-U}{E})=-\frac{t}{RC}\\[2mm]
    &\frac{E-U}{E}=e^{-\frac{t}{RC}}\\[2mm]
    &E-U=Ee^{-\frac{t}{RC}}\\[2mm]
    &U=E(1-e^{\frac{t}{RC}})
\end{align*}
这就得到了电容在充电过程中，电压关于时间的函数
\begin{BoxEquation}[电容充电过程的电压]
    U=E(1-e^{\frac{t}{RC}})
\end{BoxEquation}
\Refx{公式}{电容充电过程的电压}中，当$t\rightarrow \infty$时，$e^{-\frac{t}{RC}}\rightarrow 0$，$U\rightarrow E$。

\Refx{公式}{电容充电过程的电压}的函数图像如\Refx{图}{电容器的充电}所示，即电容在充电时，电压最初上升较快，充电速度随着电压的升高而减缓，最终趋近于电源电动势。

同时，电容和电阻越小充电速度就越快，电容和电阻越大充电速度就越慢。
\begin{Figure}{电容器的充电与放电}
    \subfigure[电容器的充电]
    {
        \label{图:电容器的充电}
        \inputfigure[scale=0.75]{Chapter09B_Figure04}
    }\qquad\quad
    \subfigure[电容器的放电]
    {
        \label{图:电容器的放电}
        \inputfigure[scale=0.75]{Chapter09B_Figure05}
    }
\end{Figure}
由此可见，电容器的充电并不是瞬间完成的，而是一个需要时间过程。电容器的电压和电荷量也不会随着充电过程的进行而无限制增大，两者分别趋向于$E$和$CE$。

\subsubsection{电容的放电过程}
电容的放电过程与充电过程是相反的，电容放电时可以视为一个电压不断降低的电源，同时假定电容在先前的充电过程中是“充满电的”，即放电过程中电容的初始电压为$E$。

现在我们要求解电容器的电压$U$关于时间的函数，列出以下方程组
\begin{Equation}[电容放电1]
    \begin{cases}
        U=IR\\[2mm]
        I=-\Fdif{Q}{t}\\[2mm]
        Q=CU
    \end{cases}
\end{Equation}
对式\ref{式:电容放电1}中的(2)做一些说明，之所以此处的导数前具有一个负号，是因为在电容的放电过程中，电流以正方向由电容的正级流经电阻回到负级时，电容上的电荷在不断被中和，是不断减少的，因此这里电流$I$是电荷量$Q$的负变化率，从而导数前有一个额外的负号。

将式\ref{式:电容放电1}中的(2)式和(3)式先后代入(1)式，消去$I,Q$
\begin{Equation}[电容放电2]
    RC\Fdif{U}{t}+U=0
\end{Equation}
式\ref{式:电容放电2}可以通过分离变量法求解，分离变量并两侧分别积分得
\begin{Equation}[电容放电3]
    \Int[E][U]{\frac{RC}{-U}\dif{U}}=
    \Int[0][t]{\dif{t}}
\end{Equation}
式\ref{式:电容放电3}中的积分下限为$E$，这是因为电容在放电时的初始电压为$E$。

现在求解式\ref{式:电容放电3}
\begin{align*}
    -&RC\ln(U)|_E^U=t\\[2mm]
    -&RC[\ln U-\ln E]=t\\[2mm]
    &\ln(\frac{U}{E})=-\frac{t}{RC}\\[2mm]
    &\frac{U}{E}=e^{-\frac{t}{RC}}\\[2mm]
    &U=Ee^{-\frac{t}{RC}}
\end{align*}
这就得到了电容在放电过程中，电压关于时间的函数
\begin{BoxEquation}[电容放电过程的电压]
    U=Ee^{-\frac{t}{RC}}
\end{BoxEquation}
\Refx{公式}{电容放电过程的电压}中，当$t\rightarrow \infty$时，$e^{-\frac{t}{RC}}\rightarrow 0$，$U\rightarrow 0$。

\Refx{公式}{电容放电过程的电压}的函数图像如\Refx{图}{电容器的放电}所示，即电容在放电时，电压最初上升较快，放电速度随着电压的降低而减缓，最终趋近于零。

同时，电容和电阻越小放电速度就越快，电容和电阻越大放电速度就越慢。

\subsection{电容器的串联与并联}
我们知道电阻是可以串联或并联的，本小节将研究电容串联或并联时的特性。
\begin{Figure}{电容器的串联与并联}
    \subfigure[电容器的串联]
    {
        \label{图:电容器的串联}
        \inputfigure{Chapter09B_Figure06}
    }\qquad
    \subfigure[电容器的并联]
    {
        \label{图:电容器的并联}
        \inputfigure{Chapter09B_Figure07}
    }
\end{Figure}
当$n$个电容$C_1,C_2,\cdots,C_n$串联时，如\Refx{图}{电容器的串联}所示，不妨设电容$C_1$的正极板带电$+Q$，那么通过静电感应$C_1$的负极板带电$-Q$。但是由于$C_1$的负极板和$C_2$的正极板通过导线相连，相当于同一导体的两个部分，适用电荷守恒定律，因此电容$C_2$的正极板也将带电$+Q$，依次类推可知，\textbf{当电容器串联时，各个电容上的电荷量相同}。

根据\Refx{定义}{电容}，每个电容器上的电压为
\begin{Equation}
    U_1=\frac{Q}{C_1}\qquad U_2=\frac{Q}{C_2}\quad\cdots\quad U_n=\frac{Q}{C_n}
\end{Equation}
这就表明，\textbf{当电容器串联时，电压依电容之反比分配在各个电容器上}。

这整个串联电容器组两端的电压等于每一个电容器极板上的电压和，即
\begin{equation*}
    U=U_1+U_2+\cdots+U_n=Q\qty(\frac{1}{C_1}+\frac{1}{C_2}+\cdots+\frac{1}{C_n})
\end{equation*}
如果将整个串联电容器组等效为一个电容，那么等效电容$C=\dfrac{Q}{U}$为
\begin{BoxPrinciple}[电容器的串联]
    电容器串联时，各个电容电荷量相同，电压依电容之反比分配。

    同时，串联后的等效总电容的倒数等于各个电容的倒数和，且总电容比任意一个电容都小。
    \begin{BoxEq}
        \frac{1}{C}=\frac{1}{C_1}+\frac{1}{C_2}+\cdots+\frac{1}{C_n}
    \end{BoxEq}
\end{BoxPrinciple}
电容的串联的等效电容公式，在数学形式上，与电阻并联的公式类似。

电容在串联时，电压是按照电容反比分配，若各个电容器的耐压也刚好与电容呈反比，那么串联后的耐压即为各个电容器的耐压之和，因此，\textbf{电容器的串联，增大耐压，减小电容。}

特别的，当两个完全相同的电容器串联时，电容变为原先的一半，耐压变为原先的一倍。

电容串联的物理实质，相当于增加了等效电容器的极板间距。

当$n$个电容$C_1,C_2,\cdots,C_n$并联时，如\Refx{图}{电容器的并联}所示，根据并联的性质可知，每一支路的电压均为$U$，因此，\textbf{当电容器并联时，各个电容上的电压相同}。

根据\Refx{定义}{电容}，每个电容器上的电荷量为
\begin{Equation}
    Q_1=C_1U\qquad Q_2=C_2U\quad\cdots\quad Q_n=C_nU
\end{Equation}
这就表明，\textbf{当电容器并联时，电荷量依电容之正比分配在各个电容器上}。

这整个并联电容器组两端的电荷量等于每一个电容器极板上的电荷量之和，即
\begin{equation*}
    Q=Q_1+Q_2+\cdots+Q_n=\qty(C_1+C_2+\cdots+C_n)U
\end{equation*}
如果将整个串联电容器组等效为一个电容，那么等效电容$C=\dfrac{Q}{U}$为
\begin{BoxPrinciple}[电容器的并联]
    电容器并联时，各个电容电压相同，电荷量依电容之正比分配。

    同时，并联后的等效总电容等于各个电容的和。
    \begin{BoxEq}
        C=C_1+C_2+\cdots+C_n
    \end{BoxEq}
\end{BoxPrinciple}
电容的并联的等效电容公式，在数学形式上，与电阻串联的公式类似。

电容在并联时，电压均相同，因此并联电容器组的耐压取决于其中耐压最低的电容器，通常并联电容器时会选取耐压相同的电容器，因此，\textbf{电容器的并联，增大电容，不改变耐压。}

特别的，当两个完全相同的电容器并联时，电容变为原先的一倍，耐压并不改变。

电容并联的物理实质，相当于增加了等效电容器的极板面积。



